\chapter{Introduction}

\section{Background}
Human papillomavirus (HPV) is one of the most common sexually transmitted viruses in the world. It spreads mainly through close skin-to-skin contact, including sexual contact, and most infected people never notice symptoms. Most infections clear on their own, but persistent infection with high-risk types can cause cancers of the cervix, anus and oropharynx, and low-risk types cause genital warts~\cite{markowitz2014}. Infections with types 16 and 18 are responsible for most cervical cancers worldwide~\cite{bruni2023}.

HPV infection is most frequent in adolescents and young adults soon after sexual debut. Vaccination works best when it is given before any exposure to the virus. For that reason, routine vaccination is recommended at 11--12 years of age, with catch-up vaccination available for young women up to the age of 26 years~\cite{markowitz2014}. University students aged 18--26 therefore sit squarely in the catch-up window, and for many of them the vaccine is still useful.

\section{Cervical Cancer in Bangladesh}
Cervical cancer is a major cause of illness and death among women in Bangladesh. GLOBOCAN 2024 estimates around 11,106 new cases and 6,557 deaths a year, making up about 14\% of all cancers diagnosed in women~\cite{iarc2024}. An earlier ICO/IARC summary put the figures at 8,268 new cases and 4,971 deaths a year and ranked cervical cancer as the second most frequent cancer among women aged 15--44 years~\cite{bruni2023}. The two sources differ, probably because of changes in data sources and estimation methods, but both show a large and avoidable burden.

Screening cannot yet close this gap. National guidelines have recommended visual inspection with acetic acid every five years for women aged 30--60 since 2018, but women are not actively invited, and screening coverage remains low~\cite{bruni2023}. Primary prevention through vaccination is therefore especially important.

\section{The National HPV Vaccination Campaign}
Bangladesh had no national HPV vaccination programme before 2023. On 2 October 2023 the government began a phased campaign with support from Gavi, UNICEF and the World Health Organization~\cite{gavi2023,unicef2023}. The first phase covered the Dhaka division, a single free dose was offered at educational institutions or designated centres, and the remaining divisions followed in October 2024~\cite{gavi2023,unbd2024}. Reports describe coverage of about 93\% of eligible girls aged 10--14 years, supported by community engagement and efforts to counter rumours and misinformation~\cite{dhakatribune2024}. Plans to include the vaccine in the routine immunisation schedule have also been reported~\cite{asianews2024}.

The campaign targeted girls in grades five to nine and out-of-school girls aged 10--14 years~\cite{dailystar2024}. Older adolescents and young women, including those at university, were not covered. Population coverage therefore remains limited, and stigma and low understanding of how HPV spreads still hold back progress~\cite{salauddin2025}. Campaign statistics count doses given. They do not show how much older students know about the vaccine, how they feel about it, or whether they would take it if it were offered.

\section{Acceptance of Vaccination}
Vaccine acceptance is a person's willingness to receive a recommended vaccine.
It ranges from having already been vaccinated through intending to be vaccinated to hesitancy and refusal.
Acceptance is shaped by awareness, knowledge, attitudes, social norms and practical considerations, including cost, availability and advice from a health provider.
Knowledge and attitudes do not always translate into uptake.
In several Bangladeshi studies, large majorities expressed willingness to be vaccinated, whereas few had actually received the vaccine~\cite{banik2022,taniya2023,chowdhury2022}.

\section{Rationale and Research Gap}
Female university students are young, are within the age range for catch-up vaccination, and will soon make health decisions for themselves and their families.
They also influence peers and relatives.
Most Bangladeshi studies of HPV have been conducted among married women, working professionals or medical students~\cite{islam2018,bhuiyan2018,chowdhury2022}.
Studies of general female university students are fewer, and those that exist focus mainly on knowledge.
Evidence on \emph{acceptance} of vaccination in this group and the factors that influence it is limited.
This question is timely because a national vaccine supply exists but does not reach university students, so any decision to extend vaccination to this group should be informed by evidence of demand.

\section{Aim and Objectives}
\subsection*{Aim}
To assess the acceptance of HPV vaccination among female university students in Dhaka, Bangladesh, and to identify the factors associated with acceptance.

\subsection*{Specific objectives}
\begin{enumerate}[leftmargin=*,itemsep=0pt]
\item To describe the socio-demographic characteristics and prior HPV awareness of the respondents.
\item To estimate the proportion of students who have been vaccinated against HPV or who intend to be vaccinated.
\item To assess students' knowledge of HPV and its vaccine, and to compare healthcare and non-healthcare students.
\item To assess students' attitudes towards HPV prevention, screening and vaccination.
\item To identify the sources through which students first learned about HPV.
\item To identify perceived barriers to vaccination, reasons for not being vaccinated, and factors that would encourage vaccination.
\item To identify factors independently associated with acceptance of HPV vaccination.
\end{enumerate}

\section{Significance of the Study}
\begin{enumerate}[leftmargin=*,itemsep=2pt]
\item \textbf{Public health.} Cervical cancer is largely preventable with vaccination, and uptake depends on whether people accept the vaccine.
\item \textbf{A group outside the campaign.} University students were not covered by the national campaign. The findings show whether demand exists among them.
\item \textbf{A gap in the evidence.} The study adds data on acceptance, not only knowledge, among students from different academic fields.
\item \textbf{Practical guidance.} The results can help universities, health authorities and educators design awareness activities and decide on catch-up vaccination.
\end{enumerate}

\section{Operational Definitions}
\begin{description}[leftmargin=0pt,style=nextline,itemsep=2pt]
\item[Acceptance] Having received at least one dose of HPV vaccine, or being ``likely'' or ``very likely'' to receive it in future.
\item[Healthcare student] A student of medicine, health sciences, pharmacy or nursing.
\item[Non-healthcare student] A student of engineering, technology, science, arts, humanities, commerce or other fields.
\item[Knowledge score] Number of correct answers to 15 true/false/not-sure statements (0--15).
\item[Attitude score] Sum of 10 five-point Likert items (10--50).
\end{description}

\section{Structure of the Report}
Chapter 2 reviews related studies. Chapter 3 describes the methods. Chapter 4 presents the results. Chapter 5 discusses them, with strengths and limitations. Chapter 6 gives conclusions and recommendations. The questionnaire and ethical approval are in the appendices.
